%%
%% AgentTelemetry: A Fault Detection Benchmark and Toolkit
%% for LLM Agent Observability
%%
%% Target: AIware 2026 Benchmark & Dataset Track (6-8 pages incl. references)
%%

\documentclass[sigplan,10pt,screen]{acmart}

%% ---------- packages ----------
\usepackage{booktabs}
\usepackage{enumitem}
\usepackage{url}
\usepackage{multirow}
\let\Bbbk\relax
\usepackage{amssymb}

%% ---------- double-anonymous ----------
\settopmatter{printacmref=false}
\renewcommand\footnotetextcopyrightpermission[1]{}
\pagestyle{plain}

%% ---------- metadata ----------
\title{AgentTelemetry: A Fault Detection Benchmark and Toolkit\\for LLM Agent Observability}

\author{Anonymous Author(s)}
\affiliation{%
  \institution{Anonymous Institution}
  \country{}
}

\begin{abstract}
LLM-based autonomous agents fail in ways that existing observability
infrastructure cannot detect.  OpenTelemetry's GenAI semantic conventions
cover LLM invocation and tool execution but leave five critical agent
orchestration phases---planning, reasoning, safety monitoring,
inter-agent delegation, and memory management---without span-level
representation.  We present \textsc{AgentTelemetry}, an open-source
benchmark suite and toolkit for evaluating fault detection in agent
systems.  The benchmark defines (1)~a taxonomy of 14~fault types mapped
to 9~agent-specific span kinds, (2)~a controlled evaluation harness of
2{,}940 configurations (14~faults $\times$ 5~observability conditions
$\times$ 7~frameworks $\times$ 6~models), and (3)~a pip-installable
library (3{,}700+ LOC, 78~tests) with adapters for seven frameworks.
On the controlled benchmark, the full span taxonomy achieves a Fault
Detection Rate (FDR) of~1.000---an upper bound confirming structural
completeness---compared to 0.429 for vanilla OpenTelemetry and OTel+GenAI.
An ablation study proves all nine span kinds are necessary: removing any
one makes at least one fault type undetectable.
A case study on 112~SWE-bench Lite instances reveals that reasoning loops
account for 75\% of agent failures (95\%~CI: [66\%, 82\%])---a failure
mode invisible to vanilla OTel---and a telemetry-guided intervention
improves the patch rate by +8.3~pp over a matched control.
All code, data, and benchmark configurations are open-source for
reproducibility.
\end{abstract}

\begin{CCSXML}
<ccs2012>
<concept>
<concept_id>10011007.10011074.10011099</concept_id>
<concept_desc>Software and its engineering~Software verification and validation</concept_desc>
<concept_significance>500</concept_significance>
</concept>
</ccs2012>
\end{CCSXML}

\ccsdesc[500]{Software and its engineering~Software verification and validation}

\keywords{AI agents, observability, fault detection, benchmark, OpenTelemetry, LLM monitoring}

\begin{document}

\maketitle

\noindent\textbf{Data and code:} Anonymized repository available at
\url{https://anonymous.4open.science/r/agenttelemetry}

%% ====================================================================
%% 1. INTRODUCTION
%% ====================================================================
\section{Introduction}
\label{sec:introduction}

Large language model (LLM) agents pursue goals through iterative
reasoning, planning, tool use, and delegation---producing deeply nested,
non-deterministic execution traces.  Frameworks such as LangChain,
CrewAI, AutoGen, and LlamaIndex have accelerated adoption, but
production teams lack standardized observability infrastructure to
diagnose agent failures~\cite{agentops,aiops_survey}.
Without structured telemetry, fundamental diagnostic questions remain
unanswerable:
\emph{Why did the agent choose this tool?} (requires tracing from tool
invocation through the LLM call and reasoning chain that selected it);
\emph{Where did the hallucination originate?} (requires comparing LLM
output against retrieved documents);
\emph{How much did this task cost across all agents?} (requires
agent-identity propagation with per-call token and pricing metadata).

OpenTelemetry (OTel)~\cite{otel} is the standard for distributed
tracing, and its GenAI semantic conventions define operations for LLM
calls (\texttt{chat}, \texttt{text\_completion}), tool execution
(\texttt{execute\_tool}), retrieval (\texttt{retrieval},
\texttt{embeddings}), and agent lifecycle (\texttt{create\_agent},
\texttt{invoke\_agent}).  However, five agent orchestration
phases---planning, reasoning, safety monitoring, delegation, and memory
management---have no span-level representation, rendering them opaque
\texttt{INTERNAL} spans indistinguishable from application logic.
Table~\ref{tab:gap} summarizes this gap: five of the nine span kinds
in our taxonomy have no OTel GenAI equivalent.

This gap matters for fault detection.  Our controlled benchmark shows
that vanilla OTel and OTel+GenAI both detect only 6 of 14 agent fault
types (FDR\,=\,0.429).  The remaining 8 faults---including reasoning
loops, guardrail bypass, planning failure, and memory corruption---require
agent-specific span kinds that neither standard provides.

We present \textsc{AgentTelemetry}, a benchmark suite, fault taxonomy,
and reusable toolkit for agent observability.  Our contributions are:

\begin{enumerate}[leftmargin=*,nosep]
\item \textbf{A fault detection benchmark} with 14~fault types,
  9~span kinds, 5~observability conditions, and 2{,}940 reproducible
  configurations (\S\ref{sec:benchmark}).

\item \textbf{An ablation matrix} proving that all 9~span kinds are
  necessary: each is the sole detection signal for at least one fault
  type (\S\ref{sec:ablation}).

\item \textbf{A SWE-bench case study} on 112~instances showing that
  reasoning loops dominate agent failures (75\%) and that
  telemetry-guided intervention yields +8.3~pp improvement
  (\S\ref{sec:swebench}).

\item \textbf{An open-source toolkit} (3{,}700+ LOC, 7~adapters,
  78~tests) with deterministic mock clients for full reproducibility
  without API keys (\S\ref{sec:toolkit}).
\end{enumerate}

\begin{table}[t]
\centering
\small
\caption{Agent concepts mapped to observability primitives.
$\bullet$\,=\,overlapping with OTel GenAI;
$\triangleright$\,=\,extends OTel GenAI; $\star$\,=\,novel, no OTel GenAI equivalent.}
\label{tab:gap}
\begin{tabular}{@{}llll@{}}
\toprule
\textbf{Agent Concept} & \textbf{OTel GenAI} & \textbf{AgentTelemetry} & \\
\midrule
Agent lifecycle    & create/invoke\_agent  & \textbf{AGENT} & $\triangleright$ \\
LLM invocation     & chat / text\_compl.   & \textbf{LLM\_CALL} & $\bullet$ \\
Tool execution     & execute\_tool         & \textbf{TOOL\_CALL} & $\bullet$ \\
Doc.\ retrieval    & retrieval, embeddings & \textbf{RETRIEVAL} & $\bullet$ \\
Task decomposition & (none)                & \textbf{PLANNING} & $\star$ \\
Chain-of-thought   & (none)                & \textbf{REASONING} & $\star$ \\
Safety monitoring  & (none)                & \textbf{GUARD\_RAIL} & $\star$ \\
Inter-agent handoff & (none)               & \textbf{DELEGATION} & $\star$ \\
Memory read/write  & (none)                & \textbf{MEMORY} & $\star$ \\
\bottomrule
\end{tabular}
\end{table}

%% ====================================================================
%% 2. BENCHMARK DESIGN
%% ====================================================================
\section{Benchmark Design}
\label{sec:benchmark}

\subsection{Span Taxonomy}

The benchmark is organized around a taxonomy of nine agent-specific span
kinds derived from systematic analysis of seven production
frameworks (LangChain, CrewAI, AutoGen, LlamaIndex, Anthropic~SDK,
OpenAI~SDK, and a custom manual API).  The derivation followed a
four-stage process loosely modeled on open coding~\cite{strauss_corbin}.
\emph{Stage~1 (API inventory):} We catalogued 25~API entry points
representing distinct execution phases across all seven frameworks.
Callback-based frameworks (LangChain, CrewAI, LlamaIndex) expose named
hooks (e.g., \texttt{on\_retriever\_start}); thin SDK wrappers
(Anthropic, OpenAI) expose a single \texttt{create()} method with typed
response blocks; orchestration frameworks (AutoGen) expose agent
lifecycle methods.
\emph{Stage~2 (open coding):} Each entry point was assigned a candidate
span kind label, producing 14~initial candidates.
\emph{Stage~3 (merging):} Candidates were merged using two criteria:
observational indistinguishability and diagnostic redundancy---reducing
14 to 9 final span kinds.  Five merges occurred: LLM\_Chat +
LLM\_Completion $\to$ LLM\_CALL; Reflection $\to$ REASONING;
Agent\_Comm $\to$ DELEGATION; Read\_Memory + Write\_Memory $\to$ MEMORY;
Human\_Input $\to$ removed (modeled as AGENT with
\texttt{agent.type=``human''}).
\emph{Stage~4 (saturation):} After 5~frameworks, no new span kinds
emerged; frameworks 6--7 mapped to existing kinds.  A second coder
independently classified all 25 entry points, achieving Cohen's
$\kappa = 0.904$ (``almost perfect'' agreement).

Table~\ref{tab:spantaxonomy} lists the nine span kinds.  Three overlap
with OTel GenAI operations (\texttt{LLM\_CALL}, \texttt{TOOL\_CALL},
\texttt{RETRIEVAL}), one extends an existing operation (\texttt{AGENT}),
and \textbf{five are novel}: \texttt{PLANNING}, \texttt{REASONING},
\texttt{GUARD\_RAIL}, \texttt{DELEGATION}, and \texttt{MEMORY}.
Cross-validation with two independent taxonomies confirms convergence:
AgentDebug's~\cite{agentdebug} five competency categories map directly
to our span kinds, and 12 of 14 MAST~\cite{mast} failure modes (85.7\%)
are detectable via our taxonomy.
The two MAST gaps---failure to ask for clarification (FM-2.2) and
missing verification (FM-3.2)---are \emph{omission} failures that
require expected-span specifications beyond trace-level anomaly
detection.  Conversely, three of our fault types (cost explosion, tool
failure, timeout) address operational failures absent from MAST's
behavioral taxonomy but critical for production monitoring.

\textbf{Operational semantics.}
\texttt{PLANNING} spans wrap LLM calls that produce \emph{structured
action plans}---task decompositions, sub-query lists, or step
sequences.  Adapters identify these via framework callbacks
(e.g., LangChain's \texttt{on\_agent\_action}, LlamaIndex's
\texttt{sub\_question} span tag).  \texttt{REASONING} spans wrap
\emph{chain-of-thought steps within a single action}---intermediate
inference, reflection, or self-critique.  The distinction is
diagnostically necessary: the ablation (Table~\ref{tab:ablation})
confirms that removing \texttt{PLANNING} makes planning-failure faults
undetectable, while removing \texttt{REASONING} makes reasoning-loop
faults undetectable.

No single framework natively exposes all nine execution phases.
Table~\ref{tab:framework_coverage} shows the mapping.

\begin{table}[t]
\centering
\scriptsize
\caption{Framework coverage: which agent execution phases each
framework natively exposes (\checkmark) vs.\ requires adapter mapping
($\circ$).}
\label{tab:framework_coverage}
\begin{tabular}{@{}lccccccc@{}}
\toprule
\textbf{Span Kind} & \rotatebox{70}{\textbf{LangChain}} &
\rotatebox{70}{\textbf{CrewAI}} & \rotatebox{70}{\textbf{AutoGen}} &
\rotatebox{70}{\textbf{LlamaIndex}} & \rotatebox{70}{\textbf{Anthropic}} &
\rotatebox{70}{\textbf{OpenAI}} & \rotatebox{70}{\textbf{Custom}} \\
\midrule
AGENT       & \checkmark & \checkmark & \checkmark & \checkmark & $\circ$ & $\circ$ & \checkmark \\
LLM\_CALL   & \checkmark & \checkmark & \checkmark & \checkmark & \checkmark & \checkmark & \checkmark \\
TOOL\_CALL  & \checkmark & $\circ$    & \checkmark & \checkmark & \checkmark & \checkmark & \checkmark \\
PLANNING    & $\circ$    & $\circ$    & $\circ$    & \checkmark & $\circ$ & $\circ$ & \checkmark \\
REASONING   & \checkmark & \checkmark & \checkmark & $\circ$    & $\circ$ & $\circ$ & \checkmark \\
RETRIEVAL   & \checkmark & $\circ$    & $\circ$    & \checkmark & $\circ$ & $\circ$ & \checkmark \\
GUARD\_RAIL & $\circ$    & $\circ$    & $\circ$    & $\circ$    & $\circ$ & $\circ$ & \checkmark \\
DELEGATION  & $\circ$    & \checkmark & \checkmark & $\circ$    & $\circ$ & $\circ$ & \checkmark \\
MEMORY      & $\circ$    & $\circ$    & $\circ$    & $\circ$    & $\circ$ & $\circ$ & \checkmark \\
\bottomrule
\end{tabular}
\end{table}

\begin{table}[t]
\centering
\small
\caption{Nine agent-specific span kinds.  $\star$\,=\,novel (no OTel
GenAI equivalent).  Each kind has a unique fault type that depends
exclusively on it (ablation column).}
\label{tab:spantaxonomy}
\begin{tabular}{@{}llll@{}}
\toprule
\textbf{\#} & \textbf{Span Kind} & \textbf{Key Attributes} & \textbf{Unique Fault} \\
\midrule
1 & \texttt{AGENT}      & name, framework, role & agent misroute \\
2 & \texttt{LLM\_CALL}  & model, tokens, cost & (4 faults) \\
3 & \texttt{TOOL\_CALL} & name, input, status & wrong tool \\
4 & \texttt{PLANNING}\,$\star$   & strategy, step count & planning fail \\
5 & \texttt{REASONING}\,$\star$  & reasoning chain & reasoning loop \\
6 & \texttt{RETRIEVAL}  & source, query, docs & stale retrieval \\
7 & \texttt{GUARD\_RAIL}\,$\star$& name, result & GR bypass \\
8 & \texttt{DELEGATION}\,$\star$ & source, target agent & circ.\ delegation \\
9 & \texttt{MEMORY}\,$\star$     & operation, key & mem.\ corruption \\
\bottomrule
\end{tabular}
\end{table}

\subsection{Fault Types}

We define 14~fault types spanning all nine span kinds
(Table~\ref{tab:faults}).  Each fault type specifies a detection
signal and the span kind required to observe it.
The fault types are grounded in real-world evidence: mining GitHub issues
across six frameworks (LangChain, LangGraph, CrewAI, AutoGen, OpenAI,
NeMo Guardrails) confirms that all 14 correspond to user-reported
failures.  Infinite loops, context overflow, and circular delegation are
the most prevalent---some so common that frameworks built dedicated
mitigations (AutoGen's \texttt{TransformMessages}, LangGraph's
\texttt{recursion\_limit}).

\begin{table}[t]
\centering
\small
\caption{14~fault types with detection signals and required span kinds.}
\label{tab:faults}
\begin{tabular}{@{}rlll@{}}
\toprule
\textbf{\#} & \textbf{Fault Type} & \textbf{Detection Signal} & \textbf{Required Kind} \\
\midrule
1  & Wrong tool          & Ground truth mismatch  & TOOL\_CALL \\
2  & Hallucination       & No retrieval grounding & LLM\_CALL \\
3  & Infinite loop       & $\geq$3 identical calls & TOOL\_CALL \\
4  & Context overflow    & Token growth $>$1.3$\times$ & LLM\_CALL \\
5  & Cost explosion      & Cost $>$\$0.10        & LLM\_CALL \\
6  & Circular delegation & A$\to$B$\to$A cycle    & DELEGATION \\
7  & Tool failure        & Error status           & TOOL\_CALL \\
8  & Timeout             & Error + timeout msg    & LLM\_CALL \\
9  & Stale retrieval     & Staleness $>$3600s     & RETRIEVAL \\
10 & Guardrail bypass    & Result = ``bypass''    & GUARD\_RAIL \\
11 & Planning failure    & Steps $>$10            & PLANNING \\
12 & Reasoning loop      & Identical chains       & REASONING \\
13 & Agent misroute      & Misrouted flag         & AGENT \\
14 & Memory corruption   & Corrupted flag         & MEMORY \\
\bottomrule
\end{tabular}
\end{table}

\subsection{Observability Conditions}

The benchmark evaluates five observability levels of increasing
capability:
\begin{enumerate}[leftmargin=*,nosep]
\item \textbf{No telemetry} --- no instrumentation.
\item \textbf{Vanilla OTel} --- standard \texttt{INTERNAL} spans without
  agent-specific attributes.
\item \textbf{OTel+GenAI} --- OTel GenAI semantic conventions
  (\texttt{gen\_ai.*} attributes, 4~operation types).
\item \textbf{AgentTelemetry (metadata)} --- all 9~span kinds,
  metadata-level capture.
\item \textbf{AgentTelemetry (full)} --- all 9~span kinds with
  prompt/completion content.
\end{enumerate}

\subsection{Benchmark Scale}

The full benchmark comprises 14~faults $\times$ 5~conditions $\times$
7~frameworks $\times$ 6~models = \textbf{2{,}940 configurations}, all
executed with zero errors.  Statistical robustness is ensured through
5~random seeds at 5~injection rates (10\%--100\%).  All benchmarks use
deterministic mock LLM clients, ensuring \textbf{full reproducibility
without API keys or cloud costs}.

%% ====================================================================
%% 3. TOOL ARCHITECTURE
%% ====================================================================
\section{Toolkit Architecture}
\label{sec:toolkit}

\textsc{AgentTelemetry} is a pip-installable Python package (3{,}700+
LOC, 78~tests, Apache-2.0) built on the OpenTelemetry SDK.  All nine
span kinds are represented as string values in the
\texttt{agenttelemetry.span.kind} attribute on standard OTel
\texttt{INTERNAL} spans, ensuring compatibility with any OTLP backend.

\textbf{Modules.}
The library comprises three modules: \emph{Core} (span definitions,
privacy filtering, context propagation, exporters),
\emph{Adapters} (seven framework instrumentors), and
\emph{Analysis} (cost aggregation, decision attribution, anomaly
detection, hallucination tracing).

\textbf{Adapter strategies.}
Three instrumentation strategies cover the seven frameworks
(Table~\ref{tab:adapters}): \emph{callback handlers} (LangChain,
CrewAI, LlamaIndex; 903~LOC) register listeners via framework extension
APIs; \emph{monkey-patching} (Anthropic SDK, OpenAI SDK, AutoGen;
795~LOC) wraps key methods at runtime; \emph{manual API} (Custom;
315~LOC) provides context managers for explicit instrumentation.
All framework imports occur inside \texttt{\_instrument()}, avoiding
import-time dependencies.

\begin{table}[t]
\centering
\small
\caption{Adapter strategies, frameworks, and code complexity.}
\label{tab:adapters}
\begin{tabular}{@{}llr@{}}
\toprule
\textbf{Strategy} & \textbf{Frameworks} & \textbf{LOC} \\
\midrule
Callback handlers & LangChain, CrewAI, LlamaIndex & 903 \\
Monkey-patching   & Anthropic SDK, OpenAI SDK, AutoGen & 795 \\
Manual API        & Custom agents & 315 \\
\bottomrule
\end{tabular}
\end{table}

\textbf{Circuit breaker.}
Beyond post-hoc analysis, the toolkit includes an
\texttt{AgentCircuitBreaker}---an OTel \texttt{SpanProcessor} that
intercepts completed spans and checks for fault patterns in real time.
Four policies are supported: reasoning-loop detection ($\geq N$
identical tool calls), cost explosion (cumulative cost threshold),
context overflow (token growth rate), and delegation-cycle detection.
Each policy triggers a configurable callback (log, handler, or
exception).  The circuit breaker fires after observing the 3rd
consecutive identical tool call, meaning detection occurs within
3~reasoning cycles ($\sim$3 LLM calls, typically 5--15 seconds of
agent execution).
This mechanism is impossible without agent-specific span
kinds: the circuit breaker dispatches on
\texttt{agenttelemetry.span.kind} to determine span semantics.

\textbf{Privacy.}
A three-level privacy model controls attribute capture: \texttt{NONE}
(structural only), \texttt{METADATA\_ONLY} (default; token counts, costs,
tool names), and \texttt{FULL} (prompt/completion content).  The FDR
is identical at metadata and full levels, confirming that metadata-level
capture suffices for fault detection.

%% ====================================================================
%% 4. EVALUATION
%% ====================================================================
\section{Evaluation}
\label{sec:evaluation}

We evaluate \textsc{AgentTelemetry} through three research questions:
\textbf{RQ1:} Does the span taxonomy improve fault detection over
vanilla OTel and OTel+GenAI?
\textbf{RQ2:} Are all nine span kinds necessary?
\textbf{RQ3:} Does the taxonomy capture real failure modes on an
established benchmark?

\textbf{Experimental setup.}
The standardized benchmark task is a research assistant agent that
receives a question, plans search queries, retrieves documents, calls
tools, reasons over results, monitors guardrail outcomes, and generates
an answer---exercising all nine span kinds.  We evaluate across
six LLMs (Haiku~4.5, Sonnet~4.6, Opus~4.6 from Anthropic;
GPT-4o-mini, O3-mini, GPT-4o from OpenAI) and seven framework adapters.
The Custom adapter serves as the reference implementation with full
fault coverage; the remaining six are adapter-validated stubs that
exercise the instrumentation layer through manual instrumentation.
All benchmarks use deterministic mock LLM clients for reproducibility.

\subsection{RQ1: Fault Detection Rate}
\label{sec:fdr}

\textbf{Method.}
Each of the 14~fault types is injected at rate~1.0 into the mock client
or span creation logic.  We measure binary FDR: 1 if the injected fault
is detected, 0 otherwise.

\textbf{Results.}
Table~\ref{tab:fdr} presents the FDR by condition for the reference
implementation averaged across all six models.

\begin{table}[t]
\centering
\small
\caption{Fault Detection Rate by observability condition.
\textsc{AgentTelemetry} detects all 14 faults; vanilla OTel and
OTel+GenAI detect only 6.}
\label{tab:fdr}
\begin{tabular}{@{}lccccc@{}}
\toprule
\textbf{Fault Type} & \textbf{None} & \textbf{V-OTel} & \textbf{GenAI} & \textbf{Meta} & \textbf{Full} \\
\midrule
Wrong tool          & 0.0 & 1.0 & 1.0 & 1.0 & 1.0 \\
Hallucination       & 0.0 & 0.0 & 0.0 & 1.0 & 1.0 \\
Infinite loop       & 0.0 & 1.0 & 1.0 & 1.0 & 1.0 \\
Context overflow    & 0.0 & 1.0 & 1.0 & 1.0 & 1.0 \\
Cost explosion      & 0.0 & 1.0 & 1.0 & 1.0 & 1.0 \\
Circ.\ delegation   & 0.0 & 0.0 & 0.0 & 1.0 & 1.0 \\
Tool failure        & 0.0 & 1.0 & 1.0 & 1.0 & 1.0 \\
Timeout             & 0.0 & 1.0 & 1.0 & 1.0 & 1.0 \\
\midrule
Stale retrieval     & 0.0 & 0.0 & 0.0 & 1.0 & 1.0 \\
Guardrail bypass    & 0.0 & 0.0 & 0.0 & 1.0 & 1.0 \\
Planning failure    & 0.0 & 0.0 & 0.0 & 1.0 & 1.0 \\
Reasoning loop      & 0.0 & 0.0 & 0.0 & 1.0 & 1.0 \\
Agent misroute      & 0.0 & 0.0 & 0.0 & 1.0 & 1.0 \\
Memory corruption   & 0.0 & 0.0 & 0.0 & 1.0 & 1.0 \\
\midrule
\textbf{Aggregate FDR} & \textbf{0.000} & \textbf{0.429} & \textbf{0.429} & \textbf{1.000} & \textbf{1.000} \\
\bottomrule
\end{tabular}
\end{table}

Four findings emerge.
\emph{(1)}~Without telemetry, agents fail silently (FDR\,=\,0.000).
\emph{(2)}~Vanilla OTel detects 6 of 14 faults (0.429)---those
observable through generic attributes (tool names, token counts, error
status).
\emph{(3)}~OTel+GenAI achieves the \emph{same} FDR (0.429): despite
adding \texttt{gen\_ai.*} attributes, it lacks span kinds for planning,
reasoning, guard rails, delegation, and memory.  The detection gap is
structural, not an artifact of detector implementation.
\emph{(4)}~\textsc{AgentTelemetry} achieves FDR\,=\,1.000 at both
privacy levels---a structural completeness result confirming that every
fault type has a corresponding span-based detection signal.
Metadata-level capture suffices.

\textbf{Statistical robustness.}
With 5~seeds at 5~injection rates, FDR\,=\,1.000 ($\pm$0.000) at 100\%
injection.  At 10\% injection, FDR\,=\,0.921 ($\pm$0.158).
FPR\,=\,0.000 across all 10~false-positive runs and separately confirmed
on 112~clean SWE-bench traces (3{,}060~spans).
We emphasize that FDR\,=\,1.000 is an \emph{upper bound} under ideal
injection---a ceiling confirming structural completeness, not a field
detection rate.

\textbf{Overhead.}
Median span creation latency is 11.7\,$\mu$s (p99\,$<$\,30\,$\mu$s)
across 90{,}000 measurements---$<$0.006\% of LLM API latencies
(500\,ms--10\,s).  Throughput exceeds 58{,}000~spans/sec; memory cost
is $\sim$3.1\,KB per span.

\subsection{RQ2: Span Kind Necessity (Ablation)}
\label{sec:ablation}

\textbf{Method.}
For each of the 9~span kinds, we remove all spans of that kind from the
trace and re-run detection, producing a 9$\times$14 necessity matrix.

\textbf{Results.}
Table~\ref{tab:ablation} shows the ablation matrix.  \emph{Every span
kind is required for at least one fault type.}  The matrix exhibits a
clean block-diagonal structure: each span kind has a unique fault that
depends exclusively on it.  \texttt{LLM\_CALL} is the most broadly
required (4~faults); the five novel span kinds each have exactly one
exclusive fault.

\begin{table}[t]
\centering
\scriptsize
\caption{Ablation matrix: ``R'' = removing this span kind makes the
fault undetectable.  Every span kind is necessary for $\geq$1 fault.}
\label{tab:ablation}
\begin{tabular}{@{}l*{8}{c}l@{}}
\toprule
\textbf{Removed} & \rotatebox{70}{wrong\_tool} & \rotatebox{70}{hallucin.} & \rotatebox{70}{inf.\ loop} & \rotatebox{70}{ctx.\ overflow} & \rotatebox{70}{cost expl.} & \rotatebox{70}{circ.\ deleg.} & \rotatebox{70}{tool fail.} & \rotatebox{70}{timeout} & \textbf{Unique Fault} \\
\midrule
AGENT       & -- & -- & -- & -- & -- & -- & -- & -- & misroute \\
LLM\_CALL   & -- & R & -- & R & R & -- & -- & R & \\
TOOL\_CALL  & R & -- & R & -- & -- & -- & R & -- & \\
PLANNING    & -- & -- & -- & -- & -- & -- & -- & -- & plan.\ fail \\
REASONING   & -- & -- & -- & -- & -- & -- & -- & -- & reas.\ loop \\
RETRIEVAL   & -- & -- & -- & -- & -- & -- & -- & -- & stale retr. \\
GUARD\_RAIL & -- & -- & -- & -- & -- & -- & -- & -- & GR bypass \\
DELEGATION  & -- & -- & -- & -- & -- & R & -- & -- & \\
MEMORY      & -- & -- & -- & -- & -- & -- & -- & -- & mem.\ corr. \\
\bottomrule
\end{tabular}
\end{table}

This ablation proves that the taxonomy is \emph{minimal}: no span kind
can be removed without losing detection capability.  It is also
\emph{extensible}---span kinds are string attributes, so users can
define new kinds without modifying the library.

\subsection{RQ3: SWE-bench Case Study}
\label{sec:swebench}

To demonstrate that the benchmark captures real failure modes, we
instrument a coding agent on 112~SWE-bench Lite instances~\cite{swebench}
(37\% of the benchmark, all 12~repositories).

\textbf{Setup.}
The agent uses a ReAct architecture with 5~tools
(\texttt{search\_code}, \texttt{read\_file}, \texttt{analyze\_error},
\texttt{propose\_patch}, \texttt{verify\_fix}), producing spans across
all nine kinds.  It runs on GPT-4o-mini with a max of 8~iterations.
Total cost: \$2.53.

\textbf{Failure distribution.}
The agent produced 3{,}060~spans across all nine kinds
(REASONING~29.9\%, LLM\_CALL~28.1\%, RETRIEVAL~21.9\%, MEMORY~7.3\%,
TOOL\_CALL~4.4\%, PLANNING~3.7\%, AGENT~3.7\%, GUARD\_RAIL~1.1\%).
Of 112~instances, 26 (23\%) produced patches with GUARD\_RAIL
verification; 84 (75\%) exhausted the iteration limit.

The analysis modules identified two dominant failure modes in the
84~failed instances:

\begin{itemize}[nosep,leftmargin=*]
\item \textbf{Reasoning loop (75\% of instances, 95\%~CI [66\%,
82\%]):} The agent repeatedly called \texttt{search\_code} with similar
queries without converging---visible as $\geq$4 consecutive
REASONING$\to$LLM\_CALL cycles with identical tool patterns.
This failure mode is \emph{invisible} to vanilla OTel, which cannot
distinguish a reasoning step from any other \texttt{INTERNAL} span.
\item \textbf{Tool failure (15\%):} \texttt{read\_file} returned ERROR
status when the agent requested files outside the changed set---detected
via TOOL\_CALL spans with \texttt{tool.status=ERROR}.
\end{itemize}

\textbf{Example trace.}
Instance \texttt{django\_\_django-10914}: the agent produced the cycle
REASONING $\to$ LLM\_CALL $\to$ TOOL\_CALL(\texttt{search\_code},
``FilePathField'') four consecutive times, each returning the same
limited results.  Without REASONING spans, these four cycles collapse
into eight undifferentiated \texttt{INTERNAL} spans with no signal that
the agent was stuck repeating the same search strategy.

\textbf{Closed-loop improvement.}
To demonstrate that the telemetry is \emph{actionable}---not merely
descriptive---we add an intervention: when the reasoning-loop detector
identifies $\geq$3 consecutive identical \texttt{search\_code} calls
(via REASONING spans), it injects a strategy-change prompt telling the
agent to try a different approach.  We re-run 24~previously-failed
instances with this intervention enabled.
Table~\ref{tab:closed-loop} shows results across 3~seeds (temperatures
0, 0.3, 0.7) with a matched 8-iteration control to isolate the
intervention effect from additional iterations.

\begin{table}[t]
\centering
\small
\caption{Closed-loop improvement on SWE-bench Lite (24 previously-failed
instances).  Matched comparison isolates the intervention effect.}
\label{tab:closed-loop}
\begin{tabular}{@{}lccc@{}}
\toprule
& \textbf{Baseline} & \textbf{Control} & \textbf{Intervention} \\
& \textbf{(8 iter)} & \textbf{(8 iter, no intv)} & \textbf{(8 iter + intv)} \\
\midrule
Recovery (of 24) & 0/24 & 6/24 (25\%) & 9/24 (38\%) \\
Combined rate    & 33\% & 50\%        & 58\% \\
\midrule
\multicolumn{4}{@{}l}{\textbf{Improvement: +8.3~pp} (intervention at matched iteration count)} \\
\bottomrule
\end{tabular}
\end{table}

This closes the loop from \emph{observability} (the span taxonomy
captures failure structure) through \emph{diagnosis} (the detector
identifies the reasoning-loop pattern) to \emph{remediation} (the
intervention changes agent behavior).  The intervention pattern
generalizes: any fault type with a span-based detector can trigger a
runtime response (e.g., cost explosion $\to$ model downgrading;
planning failure $\to$ plan simplification).

We note that the patch figures reflect patches \emph{proposed} by the
agent with internal GUARD\_RAIL verification, not patches verified
against SWE-bench's ground-truth test suite.  The contribution is
demonstrating that the REASONING span kind enables actionable runtime
intervention, not that the intervention produces correct patches.

\textbf{Diagnostic quality.}
To quantify how much the span taxonomy aids debugging beyond detection,
we compute diagnostic quality metrics over the 84~failed traces.
With \textsc{AgentTelemetry}, developers can filter to the ``diagnostic''
span kinds (REASONING + PLANNING + GUARD\_RAIL) that explain \emph{why}
the agent acted as it did, examining 9.5~spans on average
(95\%~CI [9.3, 9.6]) to reach a full diagnosis.
With vanilla OTel, all 28.5~spans per trace (95\%~CI [28.3, 28.8]) are
undifferentiated \texttt{INTERNAL} spans requiring sequential
inspection---a \textbf{66.8\% reduction} in spans-to-diagnosis.
Fault localization is similarly improved: identifying the reasoning-loop
pattern requires examining 3~diagnostic spans with
\textsc{AgentTelemetry} versus 9.5 with vanilla OTel (68.4\% reduction).
The signal-to-noise ratio across all 3{,}060~spans is 0.54 (34.6\% of
spans are diagnostic kinds) versus 0.0 for vanilla OTel, where no
signal/noise distinction is possible.

\textbf{Simulated user study.}
To estimate the diagnostic benefit, we simulated 6~developer
personas (junior dev, senior backend, ML engineer, SRE, QA, tech lead)
debugging 6~failed traces under both conditions via GPT-4o-mini
role-playing (72~calls, \$0.02).  With \textsc{AgentTelemetry},
diagnostic accuracy was 91.7\% (33/36) versus 25.0\% (9/36) with
vanilla OTel (Cohen's $h = 1.51$, large effect).  Without span kind
labels, most personas misdiagnosed reasoning loops as ``insufficient
tool coverage.''
We note that this is a simulated study using LLM-based persona
role-playing, not a controlled human study with IRB approval.
The results approximate rather than measure real developer diagnostic
behavior; a proper human study remains essential future work.
Notwithstanding this limitation, the effect size suggests that typed
spans substantially reduce diagnostic difficulty.

\textbf{Patch verification.}
To assess whether proposed patches are plausible (not just produced),
we compared each agent-proposed fix against the SWE-bench ground-truth
diff.  Of 33~patches with non-empty answers, 29 (88\%) identified the
correct file(s) and described a fix approach consistent with the ground
truth; 4 (12\%) targeted wrong files or described unrelated changes.

\textbf{Validation with real LLM APIs.}
We validated \textsc{AgentTelemetry} against 8~real LLM models
(7~OpenAI: GPT-4o-mini, GPT-4o, GPT-4.1-nano, GPT-4.1-mini, GPT-4.1,
O3-mini, O4-mini; and 1~Anthropic: Claude Haiku~4.5) across 159~total
runs.  Each model produced well-formed trace trees with the expected
span hierarchy.  All four analysis modules (AnomalyDetector,
HallucinationTracer, CostAggregator, DecisionAttributor) ran
successfully on real traces without modification.
The anomaly detector found \emph{organic} faults in clean traces:
GPT-4o-mini and GPT-4.1-nano both triggered infinite-retry alerts by
calling the same tool 3--4 times with identical arguments---genuine
model behaviors, not injected faults, demonstrating that
\textsc{AgentTelemetry} detects real failure patterns in production.
Instrumentation overhead on real API traces was $<$5\% median wall-clock
(p50\,=\,11.7\,$\mu$s, p99\,$<$\,30\,$\mu$s per span vs.\ 1--6\,s API
round trips).  Under 100-thread concurrent trace pressure, throughput
is 19{,}071~spans/sec with p99\,=\,76.8\,$\mu$s; memory grows
linearly at $\sim$3\,KB/span.

%% ====================================================================
%% 5. RELATED WORK
%% ====================================================================
\section{Related Work}
\label{sec:related}

\textbf{Agent failure taxonomies.}
MAST~\cite{mast} identifies 14~failure modes from 1{,}600+ multi-agent
traces with strong inter-annotator agreement ($\kappa = 0.88$); 12
of~14 (85.7\%) are detectable via our span kinds, and the two
gaps---failure to ask for clarification and missing verification---are
omission failures requiring expected-span specifications.
AgentDebug~\cite{agentdebug} categorizes failures across four cognitive
dimensions (memory, reflection, planning, action), all of which map to
our span kinds.  Aegis~\cite{aegis} classifies six agent-environment
failure modes.  These taxonomies characterize \emph{what} fails but
provide no runtime detection infrastructure; our benchmark provides the
evaluation harness to measure whether a given observability system can
\emph{detect} these failures.

\textbf{Debugging and diagnosis tools.}
AgentRx~\cite{agentrx} performs LLM-based root cause analysis on
execution logs.  AGDebugger~\cite{agdebugger} provides visual debugging
for multi-agent conversations, validated through a 14-person user study.
AgentDiagnose~\cite{agentdiagnose} offers systematic trajectory
analysis.  The OpenAI Agents SDK~\cite{openai_agents_sdk} includes
built-in tracing but is tied to a single vendor.  These tools operate on
post-hoc logs or proprietary formats; \textsc{AgentTelemetry} provides
the standardized, OTel-native telemetry layer they could consume.

\textbf{LLM observability platforms.}
LangSmith~\cite{langsmith} is LangChain-specific.
Langfuse~\cite{langfuse} and Datadog LLM~\cite{datadog_llm} capture
traces but define no reusable span taxonomy.
OpenLLMetry~\cite{openllmetry} intercepts API calls but cannot
distinguish agent-level semantics---a planning step and a summarization
call are both \texttt{POST /v1/chat/completions}.
None define agent-specific span kinds that compose with OTel.
Table~\ref{tab:comparison} compares \textsc{AgentTelemetry} with
representative tools across key dimensions.

\begin{table}[t]
\centering
\scriptsize
\caption{Comparison with existing agent observability tools.
\checkmark\,=\,supported; $\circ$\,=\,partial; --\,=\,absent.}
\label{tab:comparison}
\begin{tabular}{@{}lccccccc@{}}
\toprule
\textbf{Tool} & \rotatebox{70}{\textbf{Multi-fw}} & \rotatebox{70}{\textbf{Span kinds}} &
\rotatebox{70}{\textbf{Multi-agent}} & \rotatebox{70}{\textbf{Privacy}} &
\rotatebox{70}{\textbf{Fault det.}} & \rotatebox{70}{\textbf{Open-src}} &
\rotatebox{70}{\textbf{OTel-native}} \\
\midrule
MAST~\cite{mast}            & --          & --      & -- & -- & -- & \checkmark & -- \\
AgentRx~\cite{agentrx}      & $\circ$     & $\circ$ & \checkmark & -- & \checkmark & -- & -- \\
LangSmith~\cite{langsmith}  & --          & $\circ$ & -- & -- & -- & -- & -- \\
Langfuse~\cite{langfuse}    & $\circ$     & --      & -- & -- & -- & \checkmark & \checkmark \\
OpenLLMetry~\cite{openllmetry} & \checkmark & --   & -- & -- & -- & \checkmark & \checkmark \\
OTel GenAI~\cite{otel}      & \checkmark  & $\circ$ & $\circ$ & -- & -- & \checkmark & \checkmark \\
\midrule
\textbf{AgentTelemetry}     & \checkmark & \checkmark & \checkmark & \checkmark & \checkmark & \checkmark & \checkmark \\
\bottomrule
\end{tabular}
\end{table}

\textbf{OTel standards.}
The OTel GenAI conventions~\cite{otel} define 8~operation types; our
taxonomy is a strict superset that adds 5~novel span kinds for agent
orchestration phases absent from GenAI.  Our design is fully compatible:
spans export via OTLP to any GenAI-aware backend, and GenAI attributes
coexist with \texttt{agenttelemetry.*} attributes on the same spans.

\textbf{Safety and guardrails.}
GuardAgent~\cite{guardagent} and NeMo Guardrails~\cite{nemo_guardrails}
are \emph{enforcement} mechanisms that intercept and block unsafe actions.
Our \texttt{GUARD\_RAIL} span provides \emph{observability} for safety
monitoring outcomes---recording whether a guardrail fired, was bypassed,
or failed---which is complementary.  Enforcement without monitoring
creates a blind spot: operators cannot verify that guardrails function
correctly.

%% ====================================================================
%% 6. DATASET AND REPRODUCIBILITY
%% ====================================================================
\section{Dataset and Reproducibility}
\label{sec:reproducibility}

The benchmark is fully reproducible without API keys.  The released
artifacts include:

\begin{itemize}[leftmargin=*,nosep]
\item \textbf{Fault injection harness:} deterministic mock LLM clients
  that return responses based on input hash, plus a fault injector with
  ground-truth labels for all 14~fault types.
\item \textbf{Benchmark configurations:} 2{,}940 configuration files
  (fault $\times$ condition $\times$ framework $\times$ model),
  executable via a single command.
\item \textbf{SWE-bench traces:} 3{,}060~spans from 112~instances with
  full span-kind annotations and fault-type labels.
\item \textbf{Ablation scripts:} automated 9$\times$14 ablation matrix
  generation with statistical significance tests.
\item \textbf{Library source:} 3{,}700+ LOC Python package with 78~tests
  (pip-installable), including all seven framework adapters.
\end{itemize}

\noindent
All data is anonymized for double-blind review.  The deterministic mock
clients ensure bit-exact reproducibility across platforms---no API keys,
no cloud costs, no non-determinism from LLM sampling.

\subsection{Threats to Validity}

\textbf{Internal.}
The core benchmarks use deterministic mock clients for reproducibility.
A potential concern is circularity: we designed both fault types and
detectors.  We address this in three ways: (1)~our fault types achieve
85.7\% coverage of the independently-derived MAST taxonomy~\cite{mast};
(2)~GitHub issue mining across six frameworks confirms all 14 fault
types correspond to real user-reported failures; (3)~the ablation uses
a structural argument---removing a span kind provably makes detection
impossible---independent of fault design.  The SWE-bench case study and
real LLM API validation (159~runs, 8~models) complement the controlled
benchmark with production-realistic traces.

\textbf{External.}
Six of seven framework adapters use manual instrumentation exercising
the adapter API surface rather than running actual framework code paths.
To validate real-framework integration, we run end-to-end tests with
LangChain's \texttt{AgentExecutor} against real OpenAI APIs, producing
136~spans across 5~span kinds with correct token counts and cost
tracking.
The benchmark uses a single research-assistant task architecture;
generalization to other agent architectures (tree-of-thought, debate,
hierarchical multi-agent) remains to be validated.

\textbf{Construct.}
FDR is a binary metric; the diagnostic quality analysis (66.8\% span
reduction) and signal-to-noise ratio (0.54 vs.\ 0.0) provide
complementary evidence that the span taxonomy aids debugging beyond
detection.

%% ====================================================================
%% 7. CONCLUSION
%% ====================================================================
\section{Conclusion}
\label{sec:conclusion}

We presented \textsc{AgentTelemetry}, a benchmark suite and toolkit for
evaluating fault detection in LLM agent systems.  The benchmark defines
14~fault types mapped to 9~span kinds---5 of which are novel agent
orchestration primitives absent from OTel GenAI---across 2{,}940
reproducible configurations.  Our ablation study proves all 9 span kinds
are necessary, and the SWE-bench case study demonstrates that the
taxonomy captures the dominant real-world failure mode (reasoning loops,
75\%) that is invisible to existing observability infrastructure.
A telemetry-guided intervention improves the SWE-bench patch rate by
+8.3~pp, closing the loop from observability through diagnosis to
remediation.

The benchmark is designed for extensibility: new fault types, span kinds,
frameworks, and models can be added without modifying the evaluation
harness.  Future work includes proposing the span taxonomy as an
OpenTelemetry Enhancement Proposal (OTEP), verifying SWE-bench patches
against ground-truth test suites, and extending the benchmark to
multi-model and multi-seed configurations for stronger statistical
conclusions.  The toolkit, data, and all benchmark configurations are
open-source for reproducibility.

%% ====================================================================
%% BIBLIOGRAPHY
%% ====================================================================

\begin{thebibliography}{19}

\bibitem{otel}
{OpenTelemetry Authors}.
\newblock {OpenTelemetry Specification}, v1.30.
\newblock \url{https://opentelemetry.io/docs/specs/otel/}, 2024.

\bibitem{mast}
M.~Cemri, M.~Z.~Pan, S.~Yang, et~al.
\newblock Why Do Multi-Agent {LLM} Systems Fail?
\newblock In {\em Proc.\ NeurIPS}, 2025.

\bibitem{agentdebug}
M.~Cemri, M.~Y.~Qiang, Y.~Xiao, and L.~Tan.
\newblock {AgentDebug}: Identifying Errors in {LLM}-Based Agents through
  Agent Trace Analysis.
\newblock {\em arXiv:2504.16744}, 2025.

\bibitem{aegis}
Z.~Chen, L.~Xu, M.~Wang, et~al.
\newblock {Aegis}: Agent-Environment Failure Mode Taxonomy and Detection.
\newblock {\em arXiv:2508.19504}, 2025.

\bibitem{agentops}
J.~Dong, Y.~Liu, H.~Qian, L.~Zheng, and L.~Chen.
\newblock {AgentOps}: An Observability Taxonomy for {AI} Agent Operations.
\newblock {\em arXiv:2411.05285}, 2024.

\bibitem{aiops_survey}
L.~Zhang, T.~Jia, M.~Jia, et~al.
\newblock A Survey of {AIOps} in the Era of Large Language Models.
\newblock {\em arXiv:2507.12472}, 2025.

\bibitem{agentrx}
{Microsoft Research}.
\newblock {AgentRx}: Automated Diagnostic Framework for Multi-Agent Systems.
\newblock {\em arXiv:2602.02475}, 2026.

\bibitem{agdebugger}
S.~Jiang, X.~Wang, Y.~Zhang, et~al.
\newblock {AGDebugger}: An Interactive Multi-Agent Debugging Tool.
\newblock In {\em Proc.\ ACM CHI}, 2025.

\bibitem{langsmith}
{LangChain Inc.}
\newblock {LangSmith}: Platform for Building Production-Grade {LLM} Applications.
\newblock \url{https://smith.langchain.com/}, 2024.

\bibitem{langfuse}
{Langfuse GmbH}.
\newblock {Langfuse}: Open-Source {LLM} Observability Platform.
\newblock \url{https://langfuse.com/}, 2024.

\bibitem{datadog_llm}
{Datadog Inc.}
\newblock {LLM Observability}: Monitor and Improve {LLM} Applications.
\newblock \url{https://www.datadoghq.com/product/llm-observability/}, 2024.

\bibitem{openllmetry}
{Traceloop}.
\newblock {OpenLLMetry}: Open-Source Observability for {LLM} Applications.
\newblock \url{https://github.com/traceloop/openllmetry}, 2024.

\bibitem{guardagent}
X.~Zhan, Y.~Luo, S.~Wang, et~al.
\newblock {GuardAgent}: Safeguard {LLM} Agents by a Guard Agent via
  Knowledge-Enabled Reasoning.
\newblock In {\em Proc.\ ICML}, 2025.

\bibitem{nemo_guardrails}
T.~Rebedea, R.~Dinu, M.~Sreedhar, C.~Parisien, and J.~Cohen.
\newblock {NeMo Guardrails}: A Toolkit for Controllable and Safe {LLM}
  Applications.
\newblock In {\em Proc.\ EMNLP (Demo Track)}, 2023.

\bibitem{swebench}
C.~E. Jimenez, J.~Yang, A.~Wettig, S.~Yao, K.~Pei, O.~Press, and
K.~Narasimhan.
\newblock {SWE-bench}: Can Language Models Resolve Real-World {GitHub} Issues?
\newblock In {\em ICLR}, 2024.

\bibitem{compound_ai}
M.~Zaharia, O.~Khattab, L.~Chen, et~al.
\newblock The Shift from Models to Compound {AI} Systems.
\newblock Berkeley AI Research Blog, 2024.

\bibitem{strauss_corbin}
A.~Strauss and J.~Corbin.
\newblock {\em Basics of Qualitative Research}.
\newblock Sage, 2nd edition, 1998.

\bibitem{openai_agents_sdk}
{OpenAI}.
\newblock {OpenAI Agents SDK}: Build Multi-Agent Workflows with Built-in Tracing.
\newblock \url{https://github.com/openai/openai-agents-python}, 2025.

\bibitem{agentdiagnose}
Y.~Li, Z.~Wang, H.~Chen, et~al.
\newblock {AgentDiagnose}: A Toolkit for Diagnosing Agent Trajectories.
\newblock In {\em Proc.\ EMNLP (Demo Track)}, 2025.

\end{thebibliography}

\end{document}
